En el curso 9-A, los estudiantes están desarrollando una microempresa en la clase de emprendimiento. Han decidido vender helados en cono dentro del colegio y desean establecer un precio adecuado para no tener pérdidas.

Cada cono tiene forma de cono circular con radio de 2 cm y altura de 10 cm. Sobre este se coloca una bola de helado con forma aproximadamente esférica de radio 2 cm. Los estudiantes saben que una caja de helado contiene 3 litros (3000 cm$^3$) de producto.

Para tomar una decisión adecuada, necesitan determinar cuántos helados pueden producir con una caja y, a partir de esto, estimar el precio mínimo de venta por unidad para no perder dinero.

\begin{enumerate}
\item Realice una representación gráfica del helado (cono + esfera), indicando sus dimensiones.
\item Calcule el volumen total de un helado.
\item Determine cuántos helados completos se pueden obtener con una caja de 3 litros.
\item Proponga un precio mínimo por helado que permita cubrir el costo total de la caja, explicando su razonamiento.
\end{enumerate}

Sustente todos sus procedimientos utilizando argumentos matemáticos claros. No se aceptan respuestas sin justificación.